\begin{frame}{RecSys - Introduzione}
\begin{itemize}
\item Software che suggerisce all'utente elementi di interesse basandosi sulle preferenze e i comportamenti passati.
    \item Migliorano l'esperienza utente
    \item Basati su tecniche di Machine Learning
\end{itemize}
\begin{figure}[h!]
    \centering
    \begin{minipage}{0.15\textwidth}
        \centering
        \includegraphics[width=\textwidth]{images/netflix.png}
    \end{minipage}\hfill
    \begin{minipage}{0.15\textwidth}
        \centering
        \includegraphics[width=\textwidth]{images/amazon.png}
    \end{minipage}\hfill
    \begin{minipage}{0.15\textwidth}
        \centering
        \includegraphics[width=\textwidth]{images/spotify.png}
    \end{minipage}\hfill
    \begin{minipage}{0.15\textwidth}
        \centering
        \includegraphics[width=\textwidth]{images/tiktok.png}
    \end{minipage}
    \caption{Alcuni famose piattaforme che utilizzano sistemi di raccomandazione}
\end{figure}
\end{frame}

\begin{frame}{RecSys - Tipologie}
\begin{itemize}
    \item \textbf{Collaborative Filtering}: basato sulle preferenze degli utenti (user-user, item-item)
    \item \textbf{Content-based Filtering}: basato sul contenuto degli item.
    \item \textbf{Knowledge-aware}: utilizzano conoscenza esterna (es. knowledge graph)
    \item \textbf{Hybrid}: combinazione delle precedenti.
\end{itemize}
\end{frame}
